\documentclass[11pt]{amsart}
\usepackage{amsmath,amssymb,amsthm}
\usepackage[margin=1.05in]{geometry}
\raggedbottom

\numberwithin{equation}{section}
\newtheorem{theorem}{Theorem}[section]
\newtheorem{proposition}[theorem]{Proposition}
\newtheorem{lemma}[theorem]{Lemma}
\newtheorem{corollary}[theorem]{Corollary}
\theoremstyle{definition}
\newtheorem{definition}[theorem]{Definition}
\theoremstyle{remark}
\newtheorem{remark}[theorem]{Remark}
\newtheorem{example}[theorem]{Example}

\newcommand{\Z}{\mathbb{Z}}
\newcommand{\Q}{\mathbb{Q}}
\newcommand{\R}{\mathbb{R}}
\newcommand{\N}{\mathbb{N}}
\newcommand{\C}{\mathbb{C}}
\newcommand{\cA}{\mathcal{A}}
\newcommand{\Prob}{\operatorname{Prob}}
\newcommand{\id}{\mathrm{id}}

\PassOptionsToPackage{hyphens}{url}
\usepackage[hidelinks,bookmarks=false]{hyperref}

\title[Sieve primes and their phases]{The statistical mechanics of sieve-theoretic primes:\\ the critical spectrum, the log-power dichotomy,\\ and two families that are never unique}
\author{Ruqing Chen}
\address{GUT Geoservice Inc., Montreal, Canada}
\email{ruqing@hotmail.com}
\thanks{Sources, code and data for the entire series: \url{https://github.com/Ruqing1963/localized-bost-connes}; this paper is in \texttt{papers/XXIV-sieve/}.}
\date{}

\begin{document}

\begin{abstract}
We compute the thermodynamics of the classical sieve-theoretic prime families for the
localized Bost--Connes systems of \cite{Main}. The critical exponents proposed are all
correct, the open/closed dichotomy has a clean and complete explanation, and the account of
high-temperature uniqueness is wrong for two of the six families, for an elementary reason
that had not been noticed.

\emph{Critical exponents.} If $\pi_S(x)\asymp x^\theta/\log^kx$ then Abel summation gives
$\beta_c(S)=\theta$, and at $\beta=\theta$ the series behaves like $\int^x dt/(t\log^kt)$.
Hence $\beta_c=1$ for twin, Sophie Germain and Chen primes; $3/4$ for
Friedlander--Iwaniec $a^2+b^4$; $2/3$ for Heath-Brown $a^3+2b^3$; $1/2$ for Landau $n^2+1$;
and $1/c$ for Piatetski-Shapiro $\lfloor n^c\rfloor$, a continuous spectrum. All six values
are as proposed.

\emph{The dichotomy is exactly the log-power.} The critical series converges iff $k>1$, so the
low-temperature phase is closed precisely for the families defined by \emph{two}
simultaneous prime conditions --- twin, Sophie Germain, Chen, where $k=2$ --- and open for
those defined by a single condition on a polynomial or fractional sequence, where $k=1$. This
is a complete answer, and it identifies the source of closedness as the second prime
condition rather than as anything about the sieve.

\emph{Two of the six are congruence-obstructed, and are never unique --- and which two is
decided by a kernel condition, not by confinement alone.} All four congruence-trapped
families sit in a single residue class: twin and Sophie Germain primes in $2\bmod3$, Landau
and the odd Friedlander--Iwaniec primes in $1\bmod4$. But confinement obstructs only when the
occupied class lies in the \emph{kernel} of the character. For Landau and
Friedlander--Iwaniec the class is $1$, the kernel itself, so the nontrivial character mod $4$
has \emph{finite} detecting set and lies in $\Xi_\beta$ for every $\beta$: a permanent
obstruction in the sense of \cite[\S4.1]{Main}, and \textbf{for these two families the
$\mathrm{KMS}_\beta$ state is never unique, at any temperature.} For twin and Sophie Germain
primes the opposite happens: $\chi(2)=-1$ for the character mod $3$, so the detecting set is
\emph{cofinite} in $S$ and the character is detected at full strength --- the class $2$
generates $(\Z/3)^\times$, and by \cite[Lem.~5.6]{Main} the admissible classes generate
$(\Z/M\Z)^\times$ for every modulus, so no congruence obstructs these families at all.
Verified numerically to $2\cdot10^5$ (single classes as stated).

Twin, Sophie Germain, Chen, Heath-Brown and Piatetski-Shapiro primes are therefore
unobstructed, and for them the equidistribution criterion is the right one wherever the
critical series diverges.

The finding is consistent with \cite[\S5]{Main}: Theorem 5.5(5) there, with Lemma 5.6,
already established that the Sophie Germain family is confined to no proper kernel, and the
present analysis confirms it --- the mod-$3$ confinement, far from obstructing, feeds the
detecting series in full.
\end{abstract}

\maketitle

\section{The critical exponent}

The framework and notation are those of \cite{Main}, cited by DOI; no unpublished companion
paper is used.

\begin{proposition}[Abel summation]\label{prop:abel}
Let $S\subseteq P$ satisfy $\pi_S(x)\asymp x^\theta/\log^kx$ with $\theta\in(0,1]$, $k\ge0$.
Then
\begin{equation}\label{eq:abel}
\sum_{p\in S,\ p\le x}p^{-\beta}\ \asymp\ \int_2^x\frac{t^{\theta-1-\beta}}{\log^kt}\,dt ,
\end{equation}
so that $\beta_c(S)=\theta$, and at the critical value $\beta=\theta$ the right side is
$\int_2^x dt/(t\log^kt)$, which converges as $x\to\infty$ if and only if $k>1$.
\end{proposition}

\begin{proof}
The counting function has density $\asymp\theta t^{\theta-1}/\log^kt$; substitute this into
the Stieltjes integral $\sum_{p\le x}p^{-\beta}=\int_2^xt^{-\beta}\,d\pi_S(t)$.
For $\beta<\theta$ the integral is
$\asymp x^{\theta-\beta}/\log^kx\to\infty$; for $\beta>\theta$ it converges. At
$\beta=\theta$, $\int dt/(t\log^kt)=[\log^{1-k}t/(1-k)]$ for $k\ne1$ and $\log\log t$ for
$k=1$.
\end{proof}

\begin{theorem}[The spectrum]\label{thm:spectrum}
\[
\begin{array}{lccll}
S & \theta & k & \beta_c(S) & \text{phase}\\\hline
\text{twin }\{p:p+2\in P\} & 1 & 2 & 1 & \text{closed }[1,\infty)\\
\text{Sophie Germain }\{p:2p+1\in P\} & 1 & 2 & 1 & \text{closed }[1,\infty)\\
\text{Chen }\{p:p+2\in P_2\} & 1 & 2 & 1 & \text{closed }[1,\infty)\\
\text{Friedlander--Iwaniec }a^2+b^4 & 3/4 & 1 & 3/4 & \text{open }(3/4,\infty)\\
\text{Heath-Brown }a^3+2b^3 & 2/3 & 1 & 2/3 & \text{open }(2/3,\infty)\\
\text{Landau }n^2+1 & 1/2 & 1 & 1/2 & \text{open }(1/2,\infty)\\
\text{Piatetski-Shapiro }\lfloor n^c\rfloor & 1/c & 1 & 1/c & \text{open }(1/c,\infty)
\end{array}
\]
\end{theorem}

\begin{proof}
Proposition~\ref{prop:abel} with the counting functions: Brun's sieve \cite{HR} gives the upper bounds
with $k=2$ for the shifted families, and Chen's theorem \cite{Chen} the matching lower bound for Chen
primes; the Friedlander--Iwaniec \cite{FI} and Heath-Brown \cite{HB} asymptotics are theorems with $k=1$; the
Landau asymptotic with $k=1$ is conjectural (Bateman--Horn); the Piatetski-Shapiro asymptotic
with $k=1$ is a theorem for $c<243/205$ by Rivat--Sargos \cite{RS}.
\end{proof}

\begin{remark}[The dichotomy is the second prime condition]\label{rem:dichotomy}
By Proposition~\ref{prop:abel} the phase is closed exactly when $k>1$, and $k=2$ occurs
exactly for the families imposing \emph{two} simultaneous primality conditions --- $p$ and
$p+2$, $p$ and $2p+1$, $p$ and $p+2\in P_2$. The polynomial and fractional families impose
one condition and have $k=1$. So the closedness observed for Sophie Germain primes in
\cite[Thm.~5.5]{Main} is not a feature of the sieve but of the second condition, and it
transfers verbatim to twins and to Chen primes, the latter unconditionally since Chen's
theorem supplies the lower bound that the twin case still lacks.
\end{remark}

\section{Congruence confinement versus obstruction}

\begin{lemma}[Elementary congruences]\label{lem:congruence}
\begin{enumerate}
\item If $p>3$ and $p+2$ is prime then $p\equiv2\pmod3$.
\item If $p>3$ and $2p+1$ is prime then $p\equiv2\pmod3$.
\item If $p=n^2+1$ is prime and $p>2$ then $n$ is even and $p\equiv1\pmod4$.
\item Squares and fourth powers are $\equiv0,1\pmod4$, so $a^2+b^4\equiv0,1,2\pmod4$; hence
an odd prime of the form $a^2+b^4$ satisfies $p\equiv1\pmod4$.
\end{enumerate}
\end{lemma}

\begin{proof}
(1) If $p\equiv1$ then $3\mid p+2$. (2) If $p\equiv1$ then $3\mid2p+1$. (3) If $n$ is odd then
$n^2+1\equiv2\pmod4$, forcing $p=2$; so $n$ is even and $p=4m^2+1$. (4) Direct.
\end{proof}

\begin{theorem}[Permanent obstruction --- and its exact scope]\label{thm:permanent}
Let $\chi_4$ be the nontrivial character modulo $4$ and $\chi_3$ the nontrivial character
modulo $3$; for a Dirichlet character, $|1-\chi(\sigma_p)|=|1-\chi(p)|$, the Frobenius
convention of \cite[(1.1)]{Main} inverting $\chi(p)$ without changing the modulus of the
summand. Then:
\begin{enumerate}
\item For $S$ the Landau primes $n^2+1$ or the Friedlander--Iwaniec primes $a^2+b^4$, the
detecting set $D_{\chi_4}=\{p\in S:\chi_4(p)\ne1\}$ is \emph{finite}, since by
Lemma~\ref{lem:congruence}(3)--(4) all odd members lie in the class $1\bmod4$, which is
$\ker\chi_4$. Hence $\sum_{D_{\chi_4}}p^{-\beta}<\infty$ for every $\beta>0$,
$\chi_4\in\Xi_\beta(\Q,S)$ for every $\beta$, the $\mathrm{KMS}_\beta$ simplex has at
least two extreme points at every temperature, and the state is \emph{never} unique.
\item For $S$ the twin or Sophie Germain primes the same confinement has the opposite
effect: by Lemma~\ref{lem:congruence}(1)--(2) all members $p>3$ lie in the class $2\bmod3$,
and $\chi_3(2)=-1\ne1$, so $D_{\chi_3}$ is \emph{cofinite} in $S$ and
$\sum_{D_{\chi_3}}p^{-\beta}=\sum_{p\in S}p^{-\beta}+O(1)$: the character is detected at
full strength. More generally the class $2$ generates $(\Z/3)^\times$, and by
\cite[Lem.~5.6]{Main} the admissible classes of these families generate
$(\Z/M\Z)^\times$ for every modulus $M$, so \emph{no} Dirichlet character has finite
detecting set: there is no permanent congruence obstruction for twin or Sophie Germain
primes.
\end{enumerate}
\end{theorem}

\begin{proof}
A character trivial on the single occupied class has finite detecting set --- that is (1),
where the class is the kernel itself --- and a character nontrivial on it detects all but
finitely many members --- that is the first half of (2). For the last claim of (2), a
character with finite detecting set would be trivial on all but finitely many of $S$, hence
on the subgroup generated by the occupied classes, which is everything by
\cite[Lem.~5.6]{Main}. The never-uniqueness in (1) is \cite[Thm.~3.9]{Main}, uniqueness
requiring $\sum_{D_\chi}p^{-\beta}=\infty$ for every nontrivial $\chi$; the obstruction is
permanent in the sense of \cite[\S4.1]{Main}.
\end{proof}

\begin{remark}[Numerical verification]\label{rem:numver}
Computed to $2\cdot10^5$, counting $p>3$ for the mod-$3$ rows, odd $p$ for Landau, all
representable primes for Friedlander--Iwaniec, $p>3$ with $p+2\in P_2$ for Chen, positive
$a,b$ for Heath-Brown (the range of \cite{HB}), and the distinct primes
$\lfloor n^c\rfloor\le2\cdot10^5$ for Piatetski-Shapiro:
\[
\begin{array}{lrcl}
S & \#S & M & \text{coprime classes occupied}\\\hline
\text{twin (smaller member)} & 2159 & 3 & \{2\}\ \text{of}\ \{1,2\}\\
\text{Sophie Germain} & 2056 & 3 & \{2\}\ \text{of}\ \{1,2\}\\
\text{Landau }n^2+1 & 63 & 4 & \{1\}\ \text{of}\ \{1,3\}\\
\text{Friedlander--Iwaniec} & 938 & 4 & \{1\}\ \text{of}\ \{1,3\}\\
\text{Chen} & 7523 & 3,4 & \text{both classes}\\
\text{Heath-Brown} & 289 & 4,9 & \text{both mod }4;\ \{1,2,7,8\}\ \text{mod }9\\
\text{Piatetski-Shapiro }c=1.1 & 6040 & 3,4 & \text{both classes}
\end{array}
\]
The single-class rows now illustrate \emph{both} halves of Theorem~\ref{thm:permanent}:
for Landau and Friedlander--Iwaniec the occupied class is the kernel $\{1\}$
(obstruction); for twin and Sophie Germain it is $\{2\}$, on which $\chi_3$ is $-1$
(full-strength detection). For Heath-Brown modulo $9$ the set $\{1,2,7,8\}$ is not
contained in a proper subgroup of
$(\Z/9)^\times$: a character trivial on it must have $\chi(2)=1$ and hence be trivial, since
$2$ generates. So Heath-Brown is not obstructed either.
\end{remark}

\begin{corollary}[The proposed criterion is not the right one]\label{cor:criterion}
The claim that high-temperature uniqueness is equivalent to equidistribution across
admissible coprime classes is correct in one direction --- equidistribution does force
$\sum_{D_\chi}p^{-\beta}\asymp\zeta_S(\beta)=\infty$ for $\beta<\beta_c$, hence uniqueness ---
but it is not necessary: twin and Sophie Germain primes are as far from equidistributed as
possible, confined to one class, and are nevertheless unobstructed, the confining class
generating the group. What decides is not equidistribution but
\cite[Thm.~5.5(5)]{Main}: confinement to a proper \emph{kernel}. That happens for Landau
and Friedlander--Iwaniec, and for no other family here.
\end{corollary}

\begin{remark}[Consequence for \cite{Main}]\label{rem:consequence}
\cite[Thm.~5.5]{Main} establishes for the Sophie Germain family the Brun-summability, the
closed low-temperature phase, and --- in its part (5), with \cite[Lem.~5.6]{Main} --- the
uniqueness criterion in the divergent range. All of it stands, and
Theorem~\ref{thm:permanent}(2) is its confirmation from the congruence side: the mod-$3$
confinement feeds the detecting series rather than starving it. The correct picture for
Sophie Germain primes is therefore: no permanent obstruction anywhere; uniqueness wherever
the series diverges (conditional on the lower bound of \cite[Thm.~5.5(2)]{Main}); and full
symmetry breaking, $\Xi_\beta=\widehat{G_S}$, on the closed phase $[1,\infty)$ where
$\zeta_S$ converges.
\end{remark}

\section{The continuous spectrum}

\begin{theorem}[Piatetski-Shapiro]\label{thm:PS}
For $1<c<243/205$ the asymptotic $\pi_{\mathrm{PS}(c)}(x)\sim x^{1/c}/\log x$ is a theorem
(Piatetski-Shapiro for $c<12/11$, Heath-Brown for $c<755/662$, Rivat--Sargos for
$c<243/205$). Hence for those $c$ the system $\cA_{\Q,S_{\mathrm{PS}}(c)}$ has
$\beta_c=1/c\in(205/243,1)$, unconditionally, with an open low-temperature phase; and
$c\mapsto1/c$ traces a continuous family of critical temperatures.
\end{theorem}

\begin{remark}[The range is narrow]\label{rem:narrow}
$205/243=0.8436$, so the unconditional range of critical temperatures is
$\beta_c\in(0.8436,1)$, an interval of length about $0.156$. The full conjectural range
$(0,1)$ requires the asymptotic for all $c>1$, which is open. It is worth stating the
unconditional interval numerically, since ``a continuous spectrum on $(0,1)$'' overstates
what is available.
\end{remark}

\section{Assessment}

\[
\begin{array}{lll}
\text{Claim} & \text{Status} & \text{Reference}\\\hline
\beta_c=1\ \text{for twin, SG, Chen} & \text{true} & \text{Thm.~1.2}\\
\beta_c=3/4,\ 2/3,\ 1/2\ \text{for FI, HB, Landau} & \text{true} & \text{Thm.~1.2}\\
\beta_c=1/c\ \text{for Piatetski-Shapiro} & \text{true} & \text{Thm.~3.1}\\
\text{closed iff Brun-summable} & \text{true; and iff }k>1 & \text{Prop.~1.1}\\
\text{closedness comes from the sieve} & \text{no: from the 2nd prime condition} & \text{Rem.~1.3}\\
\text{uniqueness}\iff\text{equidistribution} & \text{no: kernel confinement decides} & \text{Cor.~2.4}\\
\text{Landau, FI: unique somewhere} & \textbf{false: never unique} & \text{Thm.~2.2(1)}\\
\text{twin, SG obstructed} & \textbf{no: class }2\text{ generates} & \text{Thm.~2.2(2)}\\
\text{Chen, HB, PS obstructed} & \text{no} & \text{Rem.~2.3}\\
\text{PS spectrum covers }(0,1) & \text{unconditionally only }(0.8436,1) & \text{Rem.~3.2}
\end{array}
\]

Two of the three proposed pieces are correct and one is not, and the failure is elementary
rather than deep. The critical exponents are exactly the counting exponents; the open/closed
dichotomy is exactly the power of the logarithm, hence exactly the number of simultaneous
prime conditions. Four of the six families are trapped in a single residue class by a
congruence that has nothing to do with sieves --- but confinement obstructs only when the
class sits in the kernel. That happens for Landau and Friedlander--Iwaniec, whose
$\mathrm{KMS}$ states are never unique by \cite[Thm.~3.9]{Main}; for twin and Sophie
Germain primes the same confinement is detection at full strength, and there is no
obstruction at all.

The moral is worth recording because it is the opposite of the usual one in this corpus. Here
the analytic input --- the hard sieve theorems --- is not where the difficulty lies; those
theorems give the critical exponents and the phase types cleanly. What decides the
qualitative behaviour is a congruence available to a first-year student, read correctly:
$\Xi_\beta$ measures how a set of primes sits inside residue classes \emph{relative to the
kernels of characters}. A set confined to the class $1$ is maximally obstructed no matter
how delicate its density; a set confined to a generating class is not obstructed at all.

Two questions. For the closed-phase families --- twin, Sophie Germain, Chen, all now known
unobstructed --- high-temperature uniqueness coexists with the closed phase as soon as the
divergence of the critical series below $\beta_c$ is available, which is unconditional for
Chen (Chen's theorem \cite{Chen} supplies the lower bound) and conditional for the other
two; can the Chen case be carried out in full, making it the first family in the corpus with
both a closed phase and proven high-temperature uniqueness? And second, for the two
obstructed families the permanent obstruction has order $2$; is the full $\Xi_\beta$ for,
say, the Landau primes equal to that order-two group below $\beta_c$, or is it larger?

\begin{thebibliography}{9}
\bibitem{Chen} J. R. Chen, \emph{On the representation of a larger even integer as the sum of a prime and the product of at most two primes}, Sci. Sinica 16 (1973), 157--176.
\bibitem{FI} J. Friedlander, H. Iwaniec, \emph{The polynomial $X^2+Y^4$ captures its primes}, Ann. of Math. (2) 148 (1998), 945--1040.
\bibitem{HB} D. R. Heath-Brown, \emph{Primes represented by $x^3+2y^3$}, Acta Math. 186 (2001), 1--84.
\bibitem{HR} H. Halberstam, H.-E. Richert, \emph{Sieve Methods}, LMS Monographs 4, Academic Press, 1974.
\bibitem{RS} J. Rivat, P. Sargos, \emph{Nombres premiers de la forme $\lfloor n^c\rfloor$}, Canad. J. Math. 53 (2001), 414--433.
\bibitem{Main} R. Chen, \emph{KMS state uniqueness and phase transitions in localized Bost--Connes systems}, preprint, 2026, \newline\url{https://doi.org/10.5281/zenodo.22152101}.
\end{thebibliography}

\end{document}
